% concatenated from Weight_stratified_v2_260807_w_bbl.tex
\documentclass[a4paper,11pt]{article}
%
% package declarations
%
\usepackage{graphicx}
\usepackage{amssymb}
\usepackage{amsthm}
\usepackage{subcaption}
%\usepackage[]{mdframed}
%\usepackage{arabtex,utf8}   \setcode{utf8}\usepackage{graphicx}
\usepackage{amsfonts}
\usepackage{amsmath}
\usepackage{hyperref}
\usepackage{xcolor}

%
% title page declarations
%
\title{Weight-constrained cut-down de Bruijn sequences}
\author{Chris J Mitchell and Peter R Wild\\Information Security Group, Royal Holloway, University of London}
\date{7th August 2026 (v2)}

\parskip = 1ex
\parindent = 0pt

\theoremstyle{plain}
\newtheorem{lemma}{Lemma}[section]
\newtheorem{theorem}[lemma]{Theorem}
\newtheorem{corollary}[lemma]{Corollary}
\newtheorem{result}[lemma]{Result}

\theoremstyle{definition}
\newtheorem{definition}{Definition}[section]
\newtheorem{construction}{Construction}[section]
\newtheorem{example}{Example}[section]

\theoremstyle{remark}
\newtheorem{remark}{Remark}[section]
\newtheorem{note}{Note}[section]

\begin{document}


\maketitle

\begin{abstract}
In 2011/12, Ruskey, Sawada, Stevens and Williams showed that a binary cut-down de Bruijn
sequence can be constructed containing all $n$-tuples of Hamming weight either $r$ or $r+1$ for
every possible $r$, and no others.  In this paper we examine extensions of this result that
require choosing a weight-like function applying to $n$-tuples with symbols taken from an
alphabet of arbitrary (finite) size. We consider three possible such weight functions, and in
each case establish an analogous result to that of Ruskey et al.
\end{abstract}

\section{Introduction} \label{section:introduction}

The main focus of this paper is a special class of \emph{universal sequences} (or \emph{universal
cycles}), namely --- at least for the purposes of this paper --- infinite periodic sequences with
elements from a finite set in which every $n$-tuple of a certain type occurs exactly once.
Universal sequences of a range of types have been widely studied --- see, for example, Chung et
al.\ \cite{Chung92}. According to Hollmann and van Lint \cite{Hollmann97}, the term universal
sequence in this context dates back to a 1985 paper of Lempel and Cohn \cite{Lempel85}, although
they use the term in a much more restricted sense. The more general definition we use here is
consistent with the rather abstract definition given by Chroman et al. \cite{Chroman21}.

In this paper we are interested in sequences including all $n$-tuples for which some
weight-like measure is restricted.  Such sequences have been previously considered in the
binary case by Sawada et al.~\cite{Sawada11} and Ruskey et al.~\cite{Ruskey12}, who considered
binary $n$-tuples of restricted Hamming weight. In this paper we generalise the main existence
result of Ruskey et al.\ to finite alphabets of arbitrary size. In doing so we look at three
main examples of weight-like functions, one case of which was previously considered by
Gabri\'{c} et al. \cite{Gabric20}.

\subsection{Background}

In this paper we consider periodic sequences $(s_i)$ with elements from $\mathbb{Z}_k$ for some
$k$, which we refer to as $k$-ary.  For a sequence $S = (s_i)$ and $n\geq1$ we write
$\mathbf{s}_n(i) = (s_i,s_{i+1},\ldots,s_{i+n-1})$, for an $n$-tuple corresponding to a substring
or factor of $n$ consecutive symbols occurring in the sequence at positions $i,i+1,\ldots,i+n-1$.
We say such an $n$-tuple appears in the sequence.

\begin{definition}[\cite{Alhakim24a}]
A $k$-ary \emph{$n$-window sequence $S = (s_i)$} is a periodic sequence of elements from
$\mathbb{Z}_k$ ($k>1$, $n>1$) with the property that no $n$-tuple appears more than once in a
period of the sequence, i.e.\ with the property that if $\mathbf{s}_n(i) = \mathbf{s}_n(j)$ for
some $i,j$, then $i \equiv j \pmod m$ where $m$ is the period of the sequence.  Such a sequence is
also known as a \emph{cut-down de Bruijn sequence} (see, for example, Cameron et al.\
\cite{Cameron25}), and this is the term we use here.
\end{definition}

A $k$-ary de Bruijn sequence \emph{of span $n$} is then simply an $n$-window sequence in which
every $k$-ary $n$-tuple appears once in a period, i.e.\ an $n$-window sequence of maximal period.

Following Alhakim et al.\ \cite{Alhakim24a} and many previous authors, we also introduce the
$k$-ary de Bruijn digraph. For positive integers $n$ and $k$ greater than one, let $\mathbb{Z}_k^n$
be the set of all $k^n$ tuples of length $n$ with entries from the group $\mathbb{Z}_k$ of residues
modulo $k$. The order $n$ de Bruijn digraph, $B_k(n)$, is a directed graph with $\mathbb{Z}^n_k$ as
its vertex set in which, for any two vertices $\textbf{x} = (x_0,x_1,\ldots,x_{n-1})$ and
$\textbf{y} = (y_0,y_1,\ldots,y_{n-1})$, the pair $(\textbf{x},\textbf{y})$ is an edge if and only
if $y_i = x_{i+1}$ for every $i$ ($0\leq i< n-1$). We label such an edge with the $(n+1)$-tuple
$(x_0,x_1,\ldots,x_{n-1},y_{n-1})$.

Note that we have defined two ways of specifying an edge in $B_k(n)$, namely as either a pair of
vertices $(\textbf{a},\textbf{b})$, where $\textbf{a},\textbf{b}$ are $k$-ary $n$-tuples, or as a
single $k$-ary $(n+1)$-tuple $\textbf{x}$.  We refer to an edge in $B_k(n)$ and a $k$-ary $n$-tuple
interchangeably throughout.

\begin{definition}  \label{definition:Eulerian} An \emph{Eulerian digraph} is a
strongly connected digraph, i.e.\ one in which there is a directed path between any pair of
vertices, for which every vertex has in-degree equal to out-degree.  This latter property
characterises a \emph{balanced} digraph.
\end{definition}

The name derives from the fact that there exists an Eulerian circuit, i.e.\ a directed path
visiting every edge once, in a digraph if and only if the digraph is Eulerian --- see, for example,
Corollary 6.1 of Gibbons \cite{Gibbons85} or Knuth \cite[Section 2.3.4.2, Theorem G]{Knuth97}.
Moreover, there are simple and efficient algorithms for finding Eulerian circuits --- see for
example \cite[Figure 6.5]{Gibbons85} or \cite[Section 2.3.4.2, Theorem D]{Knuth97}.

There is then a simple correspondence between a directed Eulerian circuit in $B_k(n-1)$ and a de
Bruijn sequence of span $n$; this generalises in an obvious way so that a directed circuit (or
cycle) in $B_k(n-1)$ of length $m\leq k^n$ corresponds to a cut-down de Bruijn sequence of period
$m$.

We also need the following.

\begin{definition}[Necklace --- see, for example, \cite{Ruskey92}]
A $k$-ary \emph{necklace} within the de Bruijn graph $B_k(n-1)$, written as
$[a_0,a_1,\dots,a_{m-1}]$, where $m\mid n$, consists of all distinct $n$-tuples of the form
$(a_{\overline{i}},a_{\overline{i+1}},\dots,a_{\overline{n-1}},a_0,a_1,\dots,a_{\overline{i-1}})$,
for $0\leq i<m$, i.e.\ all cyclic shifts, where $\overline{s}=s\bmod m$.
\end{definition}

De Bruijn sequences have a range of real-world applications, ranging from position location on a
linear track or rotating shaft \cite{Burns92,Burns93,Peng23,Petriu98,Yamaguchi05}, to genome
assembly \cite{Compeau11}. As observed by Cameron et al.\ \cite{Cameron25}, `for some applications
it may be more convenient to produce a cycle of arbitrary length such that there are no repeated
strings', i.e.\ a cut-down de Bruijn sequence of a particular period that derives from the
specifics of the application.

A number of authors have considered the problem of constructing cut-down de Bruijn sequences with
particular properties. Jackson et al.\ \cite[Problem 477]{Jackson09} observed that a binary
$n$-window sequence exists for every possible period $m$ satisfying $n\leq m\leq2^n$. Ruskey et
al.\ \cite{Ruskey12} showed that a cut-down de Bruijn sequence can be constructed containing all
$n$-tuples with weights in the range $[a,b]$ as long as $0\leq a<b\leq n$, and such sequences were
further studied by Li et al.\ \cite{Li23}.

Sequences containing only $n$-tuples with specified weights are a particular example of what are
sometimes known as `de Bruijn sequences with forbidden subsequences', as for example examined in
the binary case by Penne \cite{Penne10} and for arbitrary alphabets by Alhakim \cite{Alhakim22}.
Tan and Shallot \cite{Tan13}, studied the more general problem of the existence of sequences
containing only the members of a pre-specified set of $n$-tuples, where they allow such $n$-tuples
to occur multiple times in a period. These are all examples of universal sequences, as introduced
at the beginning of this paper.

\subsection{Cut down de Bruijn sequences in weight-restricted subgraphs}

As observed at the start of this paper, we are concerned with a particular class of cut-down de
Bruijn sequences, as previously considered by Sawada et al.~\cite{Sawada11} and Ruskey et al.
\cite{Ruskey12}. They showed that, if $n\geq1$ and $0\leq r<n$, the sub-graph of the de Bruijn
graph $B_2(n)$ consisting of all edges corresponding to $n$-tuples of Hamming weights $r$ and
$r+1$ is Eulerian, and hence there exists a cut down de Bruijn sequence which includes all
$n$-tuples of weights $r$ and $r+1$ (and no others).

In this note we consider to what extent a similar result holds for the $k$-ary de Bruijn graph.  Of
course, in this case there are a number of ways of computing the `weight' of an $n$-tuple, and we
therefore consider what happens for three different functions that assign a weight-like value to an
$n$-tuple.

\section{Weight functions for $k$-ary alphabets}

There are many ways of generalising the notion of Hamming weight of a binary $n$-tuple to the
$k$-ary case.  We briefly introduce some key possibilities here. The `standard' generalisation of
Hamming weight to the $k$-ary case is as follows.

\begin{definition}[Hamming weight]  \label{definition:Hamming_weight}
Suppose $k\geq2$ and $n\geq1$.  If $\mathbf{a}=(a_0,a_1,\dots,a_{n-1})$ is a $k$-ary $n$-tuple,
then the \emph{Hamming weight} $w_h(\mathbf{a})$ of $\mathbf{a}$ is equal to the number of values
$i$ for which $a_i\neq0$.
\end{definition}

It should be clear that in the binary case, i.e.\ where $k=2$, the Hamming weight corresponds to
the number of ones in the $n$-tuple, i.e.\ the expected definition.  An alternative measure for
$n$-tuples, which also corresponds to Hamming weight in the case $k=2$, is as follows.

\begin{definition}[Digit sum]  \label{definition:digit_sum}
Suppose $k\geq2$, $n\geq1$ and $\mathbf{a}=(a_0,a_1,\dots,a_{n-1})$ is a $k$-ary $n$-tuple. Then
the \emph{digit sum} (or \emph{symbol sum}) of $\mathbf{a}$ is defined as
\[ w_s(\mathbf{a}) = \sum_{i=0}^{n-1}a_i \]
where we compute the sum in $\mathbb{Z}$ and treat the values $a_i$ as integers in the range
$[0,k-1]$.
\end{definition}

\begin{remark}
Note that the digit sum is referred to simply as the \emph{weight} of an $n$-tuple in
\cite[Definition 11]{Mitchell25a}. It should be clear that the digit sum is also a natural
generalisation of Hamming weight for the binary case, although the digit sum does not have the
properties one would normally expect of a distance function.
\end{remark}

Yet another alternative way of assigning a numeric value to an $n$-tuple is as follows.

\begin{definition}[\cite{Mitchell25a}, Definition 10] \label{definition:pseudoweight}
Suppose $k\geq2$, $n\geq1$ and $\mathbf{a}=(a_0,a_1,\dots,a_{n-1})$ is a $k$-ary $n$-tuple. Define
the function $f:\mathbb{Z}_k\rightarrow\mathbb{Q}$ as follows: for any $a\in\mathbb{Z}_k$ treat $a$
as an integer in the range $[0,k-1]$ and set $f(a)=a$ if $a\neq0$ and $f(a)=k/2$ if $a=0$. Then the
\emph{pseudoweight} of $\mathbf{a}$ is defined to be the sum
\[ w_p(\mathbf{a}) = \sum_{i=0}^{n-1}f(a_i) \]
where the sum is computed in $\mathbb{Q}$.
\end{definition}

This apparently curious notion was introduced to help generate special classes of cut-down de
Bruijn sequences known as orientable sequences and negative orientable sequences (see, for example,
\cite{Mitchell26a}).

In the remainder of this paper we consider to what extent generalised versions of the Ruskey et
al.\ result \cite{Ruskey12} hold for these three weight functions.  We first need the following
preliminary definition and result that simplify arguments for all three weight functions we
examine.

\begin{definition}
Suppose $k\geq2$, $n\geq1$ and $w$ is function mapping $k$-ary $n$-tuples to $\mathbb{Q}$.  Then
$w$ is said to be \emph{permutation-agnostic} if, for any $k$-ary $n$-tuple $\mathbf{a}$,
$w(\mathbf{a})=w(p(\mathbf{a}))$ for every permutation $p$ acting on the $n$ entries in
$\mathbf{a}$.
\end{definition}

The three weight functions $w_h$, $w_s$ and $w_p$ are all permutation-agnostic since they are all
defined as a function of the elements of an $n$-tuple in a way that does not depend on the
ordering.

\begin{lemma}  \label{lemma:one_change}
Suppose $k\geq2$, $n\geq2$, $r\leq s\in \mathbb{Q}$, and $w$ is permutation-agnostic function
mapping $k$ary $n$-tuples to $\mathbb{Q}$.  Let $W_k^{r,s}(n-1)$ be the subgraph of $B_k(n-1)$
consisting of the edges corresponding to $n$-tuples $\mathbf{a}$ satisfying $r\leq
w(\mathbf{a})\leq s$. Then, if $\mathbf{a}$ and $\mathbf{b}$ are edges ($n$-tuples)) in
$W_k^{r,s}(n-1)$ that satisfy $w(\mathbf{a})=r$ and $w(\mathbf{b})=s$ and that differ in a single
position, then there are directed paths from $\mathbf{a}$ to $\mathbf{b}$ and from $\mathbf{b}$ to
$\mathbf{a}$ in $W_k^{r,s}(n-1)$.
\end{lemma}

\begin{proof}
Suppose $\mathbf{a}=(a_0,a_1\dots,a_{n-1})$ and $\mathbf{b}=(b_0,b_1\dots,b_{n-1})$, where
$a_i=b_i$ for every $i$ except $i=s$ for some $s$ ($0\leq s\leq n-1$).  Since $w$ is
permutation-agnostic, if $\mathbf{a}'$ is a cyclic shift of $\mathbf{a}$ then
$w(\mathbf{a})=w(\mathbf{a}')$.  Hence there is a directed path from $\mathbf{a}$ to
$(a_s,a_{s+1},\dots,a_{n-1},a_0,a_1,\dots,a_{s-1})$ in $W_k^{r,s}(n-1)$. Clearly
$(a_s,a_{s+1},\dots,a_{n-1},a_0,a_1,\dots,a_{s-1})$ is an incoming edge to the vertex
$(a_{s+1},\dots,a_{n-1},a_0,a_1,\dots,a_{s-1})$ and
$(a_{s+1},a_{s+2},\dots,a_{n-1},a_0,a_1,\dots,a_{s-1},b_s)=(b_{s+1},b_{s+2},\dots,b_{n-1},b_0,b_1,\dots,b_s)$
is an outgoing edge from this vertex.  That is, we can extend the directed path so that it goes
from $\mathbf{a}$ to a cyclic shift of $\mathbf{b}$. Using the same argument as previously,
there is a directed path from $(b_{s+1},b_{s+2},\dots,b_{n-1},b_0,b_1,\dots,b_s)$ to
$\mathbf{b}$ in $W_k^{r,s}(n-1)$. The same argument works in reverse, and the result follows.
\end{proof}

\section{Partitioning using Hamming weight}

We start by considering the case where the weight of an $n$-tuple is simply the number of non-zero
entries, i.e.\ using the function $w_h$ as specified in Definition~\ref{definition:Hamming_weight}.
We have a result very similar to that of Ruskey et al.\ \cite{Ruskey12}.  We first need the
following notation.

\begin{definition}[The Hamming weight subgraph]
Suppose $k\geq2$, $n\geq 1$ and $0\leq r\leq s\leq n$. We write $H_k^{r,s}(n-1)$ for the subgraph
of $B_k(n-1)$ consisting of the edges corresponding to $n$-tuples $\mathbf{a}$ satisfying $r\leq
w_h(\mathbf{a})\leq s$.
\end{definition}

\begin{remark}
$H_2^{r-1,r}(n-1)$ is what Ruskey et al.\ \cite{Ruskey12} refer to as $G(\mathbf{B}_{r-1}^r(n-1))$.
\end{remark}

The following two results are straightforward to establish.

\begin{lemma}  \label{lemma:H_balanced}
If $k\geq2$, $n\geq 2$ and $0\leq r\leq s\leq n$ then $H_k^{r,s}(n-1)$ is balanced.
\end{lemma}

\begin{proof}
$H_k^{r,r}(n-1)$ is balanced for any $r$, since if $(a_1,a_2,\dots,a_{n-1})$ is a vertex in
$B_k(n-1)$ and $x\in \mathbb{Z}_k$, then $(x,a_1,a_2,\dots,a_{n-1})$ is an incoming edge in
$H_k^{r,r}(n-1)$ if and only if $(a_1,a_2,\dots,a_{n-1},x)$ is an outgoing edge in
$H_k^{r,r}(n-1)$.  But the set of edges in $H_k^{r,s}(n-1)$ is simply the union of the sets of
edges in $H_k^{t,t}(n-1)$ for every $t$ satisfying $r\leq t\leq s$, and the result follows.
\end{proof}

\begin{lemma}  \label{lemma:H_successor rule}
Suppose $k\geq2$, $n\geq 2$, $0\leq r\leq n$ and $(a_1,a_2,\dots,a_{n-1})$ is a vertex in
$H_k^{r,r}(n-1)$.  If $(a_0,a_1,\dots,a_{n-1})$ is an incoming edge to this vertex then
$(a_1,a_2,\dots,a_{n})$ is an outgoing edge if and only if either $a_0=a_n=0$ or $a_0$ and $a_n$
are both non-zero.
\end{lemma}

\begin{proof}
This follows since all edges in $H_k^{r,r}(n-1)$ must correspond to $n$-tuples with precisely $r$
non-zero entries.
\end{proof}

Before giving the main result we need the following key lemma.

\begin{lemma}  \label{lemma:H_prelims}
Suppose $n\geq2$, $k\geq2$ and $0\leq r<n$.
\begin{itemize}
\item[(i)] If $\mathbf{a}$ is an edge (i.e.\ an $n$-tuple) in $H_k^{r,r+1}(n-1)$, then there is
    a directed path in $H_k^{r,r+1}(n-1)$ from $\mathbf{a}$ to an edge $\mathbf{a}'$ satisfying
    $w_h(\mathbf{a}')=r$.
\item[(ii)] If $\mathbf{b}$ is an edge (i.e.\ an $n$-tuple) in $H_k^{r,r+1}(n-1)$, then there
    is a directed path in $H_k^{r,r+1}(n-1)$ from an edge $\mathbf{b}'$ satisfying
    $w_h(\mathbf{b}')=r+1$ to $\mathbf{b}$.
\item[(iii)] If $\mathbf{a}=(a_0,a_1,\dots,a_{n-1})$ is an edge (i.e.\ an $n$-tuple) in
    $H_k^{r,r+1}(n-1)$ such that $w_h(\mathbf{a})=r$ and $a_i=0$ for some $i$, then there is a
    directed path in $H_k^{r,r+1}(n-1)$ from $\mathbf{a}$ to
\[ (a_0,a_1,\dots,a_{i-1},a_{i+1}, a_i,a_{i+2},\dots,a_{n-1}), \]
i.e.\ $\mathbf{a}$ modified so that $a_i$ and $a_{i+1}$ are interchanged.
\item[(iv)] If $\mathbf{b}=(b_0,b_1,\dots,b_{n-1})$ is an edge (i.e.\ an $n$-tuple) in
    $H_k^{r,r+1}(n-1)$ such that $w_h(\mathbf{b})=r+1$ and $b_i=0$ for some $i$, then there is
    a directed path in $H_k^{r,r+1}(n-1)$ from
\[ (b_0,b_1,\dots,b_{i-1},b_{i+1},b_{i},b_{i+2},\dots,b_{n-1}) \]
i.e.\ $\mathbf{b}$ modified so that $b_i$ and $b_{i+1}$ are interchanged, to $\mathbf{b}$.
\end{itemize}
\end{lemma}

\begin{proof}
\begin{itemize}
\item[(i)] If $\mathbf{a}=(a_0,a_1,\dots,a_{n-1})$ satisfies $w_h(\mathbf{a})=r+1$, then
    suppose $a_i\neq0$ for some $i$ (which must exist since $w_h(\mathbf{a})>0$).  Now clearly
    there is a directed path in $H_k^{r,r+1}(n-1)$ from $\mathbf{a}$ to
    $(a_i,a_{i+1},\dots,a_{n-1},a_0,a_1,\dots,a_{i-1})$, simply by taking successive cyclic
    shifts. Because $w_h(\mathbf{a})=r+1$ and $a_i\neq0$, the edge
    $(a_{i+1},\dots,a_{n-1},a_0,a_1,\dots,a_{i-1},0)$ is an edge in $H_k^{r,r+1}(n-1)$ with
    weight $r$, and so a directed path exists in $H_k^{r,r+1}(n-1)$ from $\mathbf{a}$ to an
    edge with weight $r$. Alternatively, if $w_h(\mathbf{a})=r$, then simply set
    $\mathbf{a}=\mathbf{a}'$. Either way the result holds.

\item[(ii)] If $\mathbf{b}=(b_0,b_1,\dots,b_{n-1})$ satisfies $w_h(\mathbf{b})=r<n$, then there
    is at least one value $b_i$ such that $b_i=0$.  As in the previous case, there is a
    directed path in $H_k^{r,r+1}(n-1)$ from
    $(b_i,b_{i+1},\dots,b_{n-1},b_0,b_1,\dots,b_{i-1})$ to $\mathbf{b}$. Because
    $w_h(\mathbf{b})=r$, the edge $(1,b_{i+1},\dots,b_{n-1},b_0,b_1,\dots,b_{i-1})$ is an edge
    in $H_k^{r,r+1}(n-1)$ with Hamming weight $r+1$, and the desired directed path exists.
    Alternatively, if $w_h(\mathbf{b})=r+1$, then simply set $\mathbf{b}=\mathbf{b}'$.  Either
    way the result holds.

\item[(iii)] There is clearly a directed path from $\mathbf{a}$ to
    $(a_i,a_{i+1},\dots,a_{n-1},a_0,a_1,\dots,a_{i-1})$ in $H_k^{r,r+1}(n-1)$, simply by taking
    successive cyclic shifts. Since $\mathbf{a}$ has weight $r$ and $a_i=0$, the $n$-tuple
    \[(a_{i+1},a_{i+2},\dots,a_{n-1},a_0,a_1,\dots,a_{i-1},a_{i+1})\]
    has weight $r$ or $r+1$ and hence is in $H_k^{r,r+1}(n-1)$.  Clearly
    \[(a_{i+2},a_{i+3},\dots,a_{n-1},a_0,a_1,\dots,a_{i-1},a_{i+1},a_i)\]
    has weight the same as $\mathbf{a}$ and hence is also in $H_k^{r,r+1}(n-1)$. That is, the
    desired directed path exists.

\item[(iv)] There is a directed path from
    $(b_{i+2},b_{i+3},\dots,b_{n-1},b_0,b_1,\dots,b_{i+1})$ to $\mathbf{b}$ in
    $H_k^{r,r+1}(n-1)$, taking cyclic shifts.  Because $\mathbf{b}$ has weight $r+1$ and
    $b_i=0$, the $n$-tuple
    \[(b_i,b_{i+2},\dots,b_{n-1},b_0,b_1,\dots,b_i)\]
    has weight $r$ or $r+1$ and hence is in $H_k^{r,r+1}(n-1)$.  Also
    \[(b_{i+1},b_i,b_{i+2},\dots,b_{n-1},b_0,b_1,\dots,b_{i-1})\]
    has weight the same as $\mathbf{b}$ and so is likewise in $H_k^{r,r+1}(n-1)$.  This
    completes the proof.
\end{itemize}
\end{proof}

We can now give our main result.

\begin{theorem}  \label{theorem:Hamming_characterisation}
Suppose $k\geq2$ and $n\geq 2$. Then:
\begin{itemize}
\item[(i)] $H_k^{r,r}(n-1)$ is Eulerian if and only if $r\in\{0,1,n-1,n\}$;
\item[(ii)] $H_k^{r,r+1}(n-1)$ is Eulerian for every $r$ ($0\leq r<n$).
\end{itemize}
\end{theorem}

\begin{remark}
For the case $k=2$ this is precisely \cite[Theorem 2.4]{Ruskey12}.
\end{remark}

Observe that in the remainder of the paper we write $x^m$ as a shorthand for $m$ consecutive
repetitions of the symbol $x$.

\begin{proof}
Because of Lemma~\ref{lemma:H_balanced}, establishing the Eulerian property just requires us to
show that the subgraphs are strongly connected.
\begin{itemize}
\item[(i)] The result is trivially true for $r=0$, since there is a single directed edge with
    Hamming weight 0, the loop to and from the vertex $0^{n-1}$.  If $r=n$, then the edges in
    $H_k^{n,n}(n-1)$ are all the $n$-tuples that do not contain a zero. But this subgraph is
    isomorphic to the full de Bruijn graph $B_k(n-2)$, and hence it is clearly strongly
    connected.

Now suppose $r=1$.  The edges of $H_k^{1,1}(n-1)$ correspond to all $n$-tuples containing a
single non-zero entry.  Suppose $\mathbf{a}=(a_0,a_1,\dots,a_{n-1})$ and
$\mathbf{b}=(b_0,b_1,\dots,b_{n-1})$ are two such $n$-tuples, where $a_u\neq0$ and $b_v\neq0$
and all the other $a_i$ and $b_i$ values are zero. Then there is clearly a directed path in
$H_k^{1,1}(n-1)$ from $\mathbf{a}$ to the $n$-tuple $(a_u,0,0,\dots,0)$ involving taking
successive cyclic shifts of $\mathbf{a}$, and similarly there is also a directed path from the
$n$-tuple $(0,0,\dots,0,b_v)$ to $\mathbf{b}$. Also, $(a_u,0,0,\dots,0)$ is an incoming edge to
the vertex $0^{n-1}$ and $(0,0,\dots,0,b_v)$ is an outgoing edge from the same vertex. Hence
there is a directed path from $\mathbf{a}$ to $\mathbf{b}$ and we have established strong
connectivity.

Finally, a somewhat similar argument applies when $r=n-1$.  Suppose
$\mathbf{a}=(a_0,a_1,\dots,a_{n-1})$ and $\mathbf{b}=(b_0,b_1,\dots,b_{n-1})$ are two
$n$-tuples corresponding to edges in $H_k^{n-1,n-1}(n-1)$, where $a_u=b_v=0$ and all the other
$a_i$ and $b_i$ values are non-zero.  Then there are directed paths in $H_k^{n-1,n-1}(n-1)$:
\begin{itemize}
\item from $\mathbf{a}$ to the $n$-tuple
    $(a_{u+1},a_{u+2},\dots,a_{n-1},a_0,a_1,\dots,a_{u-1},0)$, just by taking cyclic
    shifts;
\item then from $(a_{u+1},a_{u+2},\dots,a_{n-1},a_0,a_1,\dots,a_{u-1},0)$ to
    \\$(0,b_{v+1},b_{v+2},\dots,b_{n-1},b_0,b_1,\dots,b_{v-1})$ (from repeated application of
    Lemma~\ref{lemma:H_successor rule}); and
\item finally from $(0,b_{v+1},b_{v+2},\dots,b_{n-1},b_0,b_1,\dots,b_{v-1})$ to
    $\mathbf{b}$, again taking cyclic shifts;
\end{itemize}
yielding the required directed path from $\mathbf{a}$ to $\mathbf{b}$.

Finally suppose $1<r<n-1$, and thus $n\geq4$ --- we need to show that $H_k^{r,r}(n-1)$ is not
strongly connected.  Regardless of $r$ there will exist two edges in $H_k^{r,r}(n-1)$
corresponding to the $n$-tuples $\mathbf{a}=(1^r,0^{n-r})$ and
$\mathbf{b}=(1,0,1^{r-1},0^{n-r-1})$, and all the constituent strings are non-empty since
$n\geq 4$ and $1<r<n-1$.  We claim that there is no path in $H_k^{r,r}(n-1)$ linking
$\mathbf{a}$ to $\mathbf{b}$.

This holds because, from Lemma~\ref{lemma:H_successor rule}, any edge on a path in
$H_k^{r,r}(n-1)$ starting from $\mathbf{a}$ must contain $r$ consecutive non-zero entries
(working cyclically). Similarly, any edges on a path starting from $\mathbf{b}$ must contain a
zero followed by $r-1$ consecutive non-zero entries and then another zero (again working
cyclically). That is $\mathbf{b}$ cannot be on a path starting at $\mathbf{a}$, and vice versa.

\item[(ii)] Suppose $\mathbf{a}$ and $\mathbf{b}$ are vertices in $H_k^{r,r+1}(n-1)$.  Then, by
    Lemma~\ref{lemma:H_prelims}(i) and (ii), there are directed paths in $H_k^{r,r+1}(n-1)$
    from $\mathbf{a}$ to $\mathbf{a}'$ and from $\mathbf{b}'$ to $\mathbf{b}$, where
    $\mathbf{a}'$ has weight $r$ and $\mathbf{b}'$ has weight $r+1$.  So to establish the
    result we simply need to show there is a path in $H_k^{r,r+1}(n-1)$ from any $n$-tuple with
    weight $r$ to any $n$-tuple with weight $r+1$.

Suppose $\mathbf{a}$ has weight $r$ and $\mathbf{b}$ has weight $r+1$, where $0\leq r<n$. Then
we claim there are directed paths in $H_k^{r,r+1}(n-1)$ from $\mathbf{a}$ to $(1^r,0^{n-r})$,
and from $(1^{r+1},0^{n-r-1})$ to $\mathbf{b}$. The first assertion follows first from repeated
application of Lemma~\ref{lemma:H_prelims}(iii) to $\mathbf{a}$ to give a path from
$\mathbf{a}$ to an edge whose first $r$ entries are non-zero and the last $n-r$ entries are
zero.  A path from this vertex to $(1^r,0^{n-r})$ exists from Lemma~\ref{lemma:H_successor
rule}.

An exactly analogous argument using repeated applications of
Lemmas~\ref{lemma:H_prelims}(iv) and \ref{lemma:H_successor rule} shows that there is a
directed path from $(1^{r+1},0^{n-r-1})$ to $\mathbf{b}$.  That is, to complete the proof
it remains to show that there is a directed path in $H_k^{r,r+1}(n-1)$ from
$(1^{r},0^{n-r})$ to $(1^{r+1},0^{n-r-1})$; but this is immediate since they differ in only
one digit and hence the result follows from Lemma~\ref{lemma:one_change}.

\end{itemize}
\end{proof}

\begin{remark}
A simpler somewhat less formal approach to proving part (i) of the above theorem is to observe that
the $n$-tuples corresponding to edges in $H_k^{r,r}(n-1)$ all have the same `pattern' of
occurrences of zeros, up to cyclic shift.  This follows from Lemma~\ref{lemma:H_successor rule}.
\end{remark}

Theorem~\ref{theorem:Hamming_characterisation} has the following immediate corollary.

\begin{corollary}  \label{corollary:H_characterisation}
Suppose $k\geq2$, $n\geq 2$ and $0\leq r<s\leq n$. Then $H_k^{r,s}(n-1)$ is Eulerian.
\end{corollary}

\begin{proof}
The fact that $H_k^{r,s}(n-1)$ is balanced is immediate from Lemma~\ref{lemma:H_balanced}. We know
that the set of edges in $H_k^{r,s}(n-1)$ is the union of the sets of edges in $H_k^{t,t+1}(n-1)$
for every $t$ satisfying $r\leq t<s$, and $H_k^{t,t+1}(n-1)$ is strongly connected from
theorem~\ref{theorem:Hamming_characterisation}(ii).  So suppose $\mathbf{a}_t$ and $\mathbf{a}_u$
are edges in $H_k^{r,s}(n-1)$ of weights $t$ and $u$ respectively.  Without loss of generality
suppose $t\leq u$.  If $t=u$ or $t=u+1$ then both $\mathbf{a}_t$ and $\mathbf{a}_u$ are in
$H_k^{t,t+1}(n-1)$, and hence there is a directed path from $\mathbf{a}_t$ to $\mathbf{a}_u$. So
suppose $u>t+1$.  Then let $\mathbf{a}_{t+1},\mathbf{a}_{t+2},\dots,\mathbf{a}_{u-1}$ be arbitrary
edges of weights $t+1,t+2,\dots,u-1$ respectively.  Then since $H_k^{i,i+1}(n-1)$ is strongly
connected for every $i$, there exist directed paths in $H_k^{r,s}(n-1)$ from $\mathbf{a}_i$ to
$\mathbf{a}_{i+1}$ for every $i$ ($u\leq i<t$), and hence from $\mathbf{a}_t$ to $\mathbf{a}_u$. It
follows that $H_k^{r,s}(n-1)$ is strongly connected.
\end{proof}

\section{Partitioning using digit sum}

The second `measure' for $n$-tuples that we consider is the digit sum, as in
Definition~\ref{definition:digit_sum}.  For this case we need the following.

\begin{definition}[The digit sum subgraph]
Suppose $k\geq2$, $n\geq 2$ and $0\leq r\leq s\leq n(k-1)$. We write $S_k^{r,s}(n-1)$ for the
subgraph of $B_k(n-1)$ consisting of the edges corresponding to $n$-tuples $\mathbf{a}$ satisfying
$r\leq w_s(\mathbf{a})\leq s$.
\end{definition}

A special case of the problem we address here has been previously considered for this
function.

\begin{result}[\cite{Gabric20}]  \label{result:old_S_lemma}
Suppose $k\geq2$ and $n\geq 2$. Then $S_k^{0,s}(n-1)$ and $S_k^{s,n(k-1)}(n-1)$ are Eulerian,
for every $s$ ($0\leq s\leq n(k-1)$).
\end{result}

In fact a much more general result can be established.  Analogously to
Lemmas~\ref{lemma:H_balanced} and \ref{lemma:H_successor rule} we first have the following.

\begin{lemma}  \label{lemma:S_balanced}
If $k\geq2$, $n\geq 2$ and $0\leq r\leq s\leq n(k-1)$ then $S_k^{r,s}(n-1)$ is balanced.
\end{lemma}

\begin{proof}
The fact that $S_k^{r,r}(n-1)$ is balanced for any $r$ is immediate, since if
$(a_1,a_2,\dots,a_{n-1})$ is a vertex in $B_k(n-1)$ and $x\in \mathbb{Z}_k$, then
$(x,a_1,a_2,\dots,a_{n-1})$ is an incoming edge in $S_k^{r,r}(n-1)$ if and only if
$(a_1,a_2,\dots,a_{n-1},x)$ is an outgoing edge in $S_k^{r,r}(n-1)$.  But the set of edges in
$S_k^{r,s}(n-1)$ is simply the union of the sets of edges in $S_k^{t,t}(n-1)$ for every $t$
satisfying $r\leq t\leq s$, and the result follows.
\end{proof}

\begin{lemma}  \label{lemma:S_successor rule}
Suppose $k\geq2$, $n\geq 2$, $0\leq r\leq n(k-1)$ and $(a_1,a_2,\dots,a_{n-1})$ is a vertex in
$S_k^{r,r}(n-1)$. If $(a_0,a_1,\dots,a_{n-1})$ is an incoming edge to this vertex then
$(a_1,a_2,\dots,a_{n})$ is an outgoing edge if and only if $a_0=a_n$.
\end{lemma}

\begin{proof}
This is immediate since all edges in $S_k^{r,r}(n-1)$ correspond to $n$-tuples with digit sum $r$.
\end{proof}

Since $w_s(\mathbf{a})=w_h(\mathbf{a})$ if $k=2$, Theorem~\ref{theorem:Hamming_characterisation}
applies in this case.  Thus we restrict our attention here to the case $k>2$.  Before giving the
main result we need the following key lemma.

\begin{lemma}  \label{lemma:S_prelims}
Suppose $n\geq2$, $k>2$ and $0\leq r<n(k-1)$.
\begin{itemize}
\item[(i)] If $\mathbf{a}$ is an edge (i.e.\ an $n$-tuple) in $S_k^{r,r+1}(n-1)$, then there is
    a directed path in $S_k^{r,r+1}(n-1)$ from $\mathbf{a}$ to an edge $\mathbf{a}'$ satisfying
    $w_s(\mathbf{a}')=r$.
\item[(ii)] If $\mathbf{b}$ is an edge (i.e.\ an $n$-tuple) in $S_k^{r,r+1}(n-1)$, then there
    is a directed path in $S_k^{r,r+1}(n-1)$ from an edge $\mathbf{b}'$ satisfying
    $w_s(\mathbf{b}')=r+1$ to $\mathbf{b}$.
\item[(iii)] If $\mathbf{a}=(a_0,a_1,\dots,a_{n-1})$ is an edge (i.e.\ an $n$-tuple) in
    $S_k^{r,r+1}(n-1)$ such that $w_s(\mathbf{a})=r$, $a_i<k-1$ and $a_{i+1}>0$, for some $i$,
    then there is a directed path in $S_k^{r,r+1}(n-1)$ from $\mathbf{a}$ to
\[ (a_0,a_1,\dots,a_{i-1},a_i+1, a_{i+1}-1,a_{i+2},\dots,a_{n-1}), \]
i.e.\ $\mathbf{a}$ modified so that $a_i$ is increased by 1 and $a_{i+1}$ decreased by 1.
\item[(iv)] If $\mathbf{b}=(b_0,b_1,\dots,b_{n-1})$ is an edge (i.e.\ an $n$-tuple) in
    $S_k^{r,r+1}(n-1)$ such that $w_s(\mathbf{b})=r+1$, $b_i<k-1$ and $b_{i+1}>0$, for some
    $i$, then there is a directed path in $S_k^{r,r+1}(n-1)$ from
\[ (b_0,b_1,\dots,b_{i-1},b_i+1,b_{i+1}-1,b_{i+2},\dots,b_{n-1}) \]
i.e.\ $\mathbf{b}$ modified so that $b_i$ is increased by 1 and $b_{i+1}$ decreased by 1, to
$\mathbf{b}$.
\end{itemize}
\end{lemma}

\begin{proof}
\begin{itemize}
\item[(i)] If $\mathbf{a}=(a_0,a_1,\dots,a_{n-1})$ satisfies $w_s(\mathbf{a})=r+1$, then
    suppose $a_i\neq0$ for some $i$ (which must exist since $w_s(\mathbf{a})>0$).  Now using
    cyclic shifting there is a directed path in $S_k^{r,r+1}(n-1)$ from $\mathbf{a}$ to
    $(a_i,a_{i+1},\dots,a_{n-1},a_0,a_1,\dots,a_{i-1})$. Because $w_s(\mathbf{a})=r+1$, the
    edge $(a_{i+1},\dots,a_{n-1},a_0,a_1,\dots,a_{i-1},a_{i}-1)$ is an edge in
    $S_k^{r,r+1}(n-1)$ with digit sum $r$, and so a directed path exists in $S_k^{r,r+1}(n-1)$
    from $\mathbf{a}$ to an edge with digit sum $r$. Alternatively, if $w_s(\mathbf{a})=r$,
    then simply set $\mathbf{a}=\mathbf{a}'$. Either way the result holds.
\item[(ii)] If $\mathbf{b}=(b_0,b_1,\dots,b_{n-1})$ satisfies $w_s(\mathbf{b})=r<n(k-1)$, then
    there is at least one value $b_i$ such that $b_i<k-1$.  As in the previous case, using
    cyclic shifting there is a directed path in $S_k^{r,r+1}(n-1)$ from
    $(b_i,b_{i+1},\dots,b_{n-1},b_0,b_1,\dots,b_{i-1})$ to $\mathbf{b}$. Because
    $w_s(\mathbf{b})=r$, the edge $(b_i+1,b_{i+1},\dots,b_{n-1},b_0,b_1,\dots,b_{i-1})$ is an
    edge in $S_k^{r,r+1}(n-1)$ with digit sum $r+1$, and the desired directed path exists.
    Alternatively, if $w_s(\mathbf{b})=r+1$, then simply set $\mathbf{b}=\mathbf{b}'$.  Either
    way the result holds.
\item[(iii)] There is clearly a directed path from $\mathbf{a}$ to
    $(a_i,a_{i+1},\dots,a_{n-1},a_0,a_1,\dots,a_{i-1})$ in $S_k^{r,r+1}(n-1)$ simply from
    cyclic shifting. Since $\mathbf{a}$ has digit sum $r$, the $n$-tuple
    \[(a_{i+1},a_{i+2},\dots,a_{n-1},a_0,a_1,\dots,a_{i-1},a_i+1)\]
    has digit sum $r+1$ and hence is in $S_k^{r,r+1}(n-1)$.  It follows that
    \[(a_{i+2},a_{i+3},\dots,a_{n-1},a_0,a_1,\dots,a_{i-1},a_i+1,a_{i+1}-1)\]
    has digit sum $r$ and hence is also in $S_k^{r,r+1}(n-1)$. Finally, there is a directed
    path from this latter $n$-tuple to
    \[ (a_0,a_1,\dots,a_{i-1},a_i+1, a_{i+1}-1,a_{i+2},\dots,a_{n-1}), \]
    via cyclic shifting and so the desired directed path exists.
\item[(iv)] There is a directed path from
    $(b_{i+2},b_{i+3},\dots,b_{n-1},b_0,b_1,\dots,b_{i+1})$ to $\mathbf{b}$ in
    $S_k^{r,r+1}(n-1)$ using cyclic shifting.  Because  $\mathbf{b}$ has digit sum $r+1$, the
    $n$-tuple
\[(b_{i+1}-1,b_{i+2},\dots,b_{n-1},b_0,b_1,\dots,b_i)\]
has digit sum $r$ and hence is in $S_k^{r,r+1}(n-1)$.  Hence
\[(b_i+1,b_{i+1}-1,b_{i+2},\dots,b_{n-1},b_0,b_1,\dots,b_{i-1})\]
has digit sum $r+1$ and so is likewise in $S_k^{r,r+1}(n-1)$. Finally, there is a directed path
    from
    \[ (b_0,b_1,\dots,b_{i-1},b_i+1, b_{i+1}-1,b_{i+2},\dots,b_{n-1}), \]
    to this latter $n$-tuple via cyclic shifting and so the desired directed path exists.  This
    completes the proof.
\end{itemize}
\end{proof}

We can now give the following result.

\begin{theorem}  \label{theorem:S_characteristic}
Suppose $n\geq2$ and $k>2$. Then:
\begin{itemize}
\item[(i)] $S_k^{r,r}(n-1)$ is Eulerian if and only if $r\in\{0,1,(k-1)n-1,(k-1)n\}$;
\item[(ii)] $S_k^{r,r+1}(n-1)$ is Eulerian for every $r$ ($0\leq r<(k-1)n$).
\end{itemize}
\end{theorem}

\begin{proof}
\begin{itemize}
\item[(i)] If $r=0$ or $r=(k-1)n$ then the only $n$-tuple with digit sum $r$ is $0^n$ or
    $(k-1)^n$, respectively.  The result follows.  Somewhat similarly, the only $n$-tuple with
    digit sum $r=1$ consists of $n-1$ zeros and a single 1 symbol.  Clearly this set of $n$
    $n$-tuples can be arranged in a single cycle --- a necklace.  A precisely analogous
    argument holds for $r=(k-1)n-1$.

    Now suppose $1<r<n-1$, and hence $n\geq4$.  By Lemma~\ref{lemma:S_successor rule}, the only
cycles in $S_k^{r,r}(n-1)$ consist of necklaces, and thus have period dividing $n$.  However,
if $1<r<(k-1)n-1$ then $S_k^{r,r}(n-1)$ will clearly contain more than $n$ edges ($n$-tuples),
and hence it contains more than one necklace.  However, two necklaces in $S_k^{r,r}(n-1)$
cannot share a vertex since, because $w_s$ is one-to-one, there is only at most one edge of
weight $r$ outgoing from (and incoming to) any vertex. The desired result follows.
\item[(ii)] Just as previously, because of Lemma~\ref{lemma:S_balanced}, establishing the
    Eulerian property just requires us to show that the subgraphs are strongly connected.
    Suppose $\mathbf{a}$ and $\mathbf{b}$ are vertices in $S_k^{r,r+1}(n-1)$.  Then, by
    Lemma~\ref{lemma:S_prelims}(i) and (ii), there are directed paths in $S_k^{r,r+1}(n-1)$
    from $\mathbf{a}$ to $\mathbf{a}'$ and from $\mathbf{b}'$ to $\mathbf{b}$, where
    $\mathbf{a}'$ has digit sum $r$ and $\mathbf{b}'$ has digit sum $r+1$.  So to establish the
    result we simply need to show there is a path in $S_k^{r,r+1}(n-1)$ from any $n$-tuple with
    digit sum $r$ to any $n$-tuple with digit sum $r+1$.

Suppose $\mathbf{a}$ has digit sum $r$ and $\mathbf{b}$ has digit sum $r+1$, where
$r=s(k-1)+t$, $s\geq 0$, and $0\leq t<k-1$. Then we claim there are directed paths in
$S_k^{r,r+1}(n-1)$ from $\mathbf{a}$ to $((k-1)^s,t,0^{n-s-1})$, and from
$((k-1)^s,t+1,0^{n-s-1})$ to $\mathbf{b}$.  The first assertion follows from repeated
application of Lemma~\ref{lemma:S_prelims}(iii) to $\mathbf{a}$ until an $n$-tuple
$(c_0,c_1,\dots,c_{n-1})$ of digit sum $r$ is reached for which there is no $i$ such that
$c_i<k-1$ and $c_{i+1}>0$.  It follows that $\mathbf{c}=((k-1)^s,t,0^{n-s-1})$, and the
first claim is established.  An exactly analogous argument using repeated applications of
Lemma~\ref{lemma:S_prelims}(iv) shows that there is a directed path from
$((k-1)^s,t+1,0^{n-s-1})$ to $\mathbf{b}$.  That is, to complete the proof it remains to
show that there is a directed path in $S_k^{r,r+1}(n-1)$ from $((k-1)^s,t,0^{n-s-1})$ to
$((k-1)^s,t+1,0^{n-s-1})$; but this is immediate since they differ in only one digit, and
hence the result follows from Lemma~\ref{lemma:one_change}.
\end{itemize}
\end{proof}

Clearly Result~\ref{result:old_S_lemma} follows from Theorem~\ref{theorem:S_characteristic} as
a special case.

Analogously to the Hamming weight case, we have the following simple corollary, whose proof is
immediate using the same arguments as for Corollary~\ref{corollary:H_characterisation}.

\begin{corollary} \label{corollary:S_characteristic}
Suppose $n\geq2$, $k>2$ and $0\leq r< s\leq (k-1)n$. Then $S_k^{r,s}(n-1)$ is Eulerian.
\end{corollary}

\section{Partitioning using pseudoweight}

The third `measure' for $n$-tuples that we consider is the pseudoweight, as in
Definition~\ref{definition:pseudoweight}.  In this case we need the following.

\begin{definition}[The pseudoweight subgraph]
Suppose $k>2$, $n\geq 2$ and $0\leq r\leq s\leq (k-1)n$ (where $r\in\mathbb{Z}$ if $k$ is even and
$r=t/2$ for some $t\in\mathbb{Z}$ if $k$ is odd). We write $P_k^{r,s}(n-1)$ for the subgraph of
$B_k(n-1)$ consisting of the edges corresponding to $n$-tuples $\mathbf{a}$ satisfying $r\leq
w_p(\mathbf{a})\leq s$.
\end{definition}

Exactly as for digit sum, a special case of the problem we address here has been previously
considered for this function. The discussion following \cite[Theorem 3.8]{Mitchell25a} establishes
the following result.

\begin{result}[\cite{Mitchell25a}]  \label{result:old_P_lemma}
Suppose $k\geq2$ and $n\geq 2$. Then $P_k^{n,s}(n-1)$ is Eulerian, where
\[ s = \begin{cases}
\frac{nk-1}{2}\text{~~if~~$k$ is odd}\\
\frac{nk-2}{2}\text{~~if~~$k$ is even}
\end{cases}\]
\end{result}

\begin{remark}
Since $\frac{nk}{2}$ is the average value of the function $w_p$ when calculated over all
$n$-tuples, this simply says that the subgraph of $B_k(n-1)$ consisting of all $n$-tuples with
pseudoweight less than the average value is Eulerian.
\end{remark}

If $k=2$ then $f(0)=f(1)=1$ and the pseudoweight function is trivial and uninteresting. Thus, in
the remainder of this discussion we restrict our attention to $k>2$.

First, analogously to Lemmas~\ref{lemma:H_balanced} and \ref{lemma:H_successor rule} we have the
following.

\begin{lemma}  \label{lemma:P_balanced}
If $k>2$, $n\geq 2$ and $n\leq r\leq s\leq n(k-1)$ (where $r,s\in\mathbb{Z}$ if $k$ is even and
$r=t/2$ and $s=u/2$ for some $t,u\in\mathbb{Z}$ if $k$ is odd) then $P_k^{r,s}(n-1)$ is balanced.
\end{lemma}

\begin{proof}
The fact that $P_k^{r,r}(n-1)$ is balanced for any $r$ is immediate since if
$(a_1,a_2,\dots,a_{n-1})$ is a vertex in $B_k(n-1)$ and $x\in \mathbb{Z}_k$, then
$(x,a_1,a_2,\dots,a_{n-1})$ is an incoming edge in $P_k^{r,r}(n-1)$ if and only if
$(a_1,a_2,\dots,a_{n-1},x)$ is an outgoing edge in $P_k^{r,r}(n-1)$, since they have the same
pseudoweight. But the set of edges in $P_k^{r,s}(n-1)$ is simply the union of the sets of edges in
$P_k^{t,t}(n-1)$ for every $t$ satisfying $r\leq t\leq s$, and the result follows.
\end{proof}

\begin{lemma}  \label{lemma:P_successor rule}
Suppose $k>2$, $n\geq 2$, $n\leq r\leq n(k-1)$ and $(a_1,a_2,\dots,a_{n-1})$ is a vertex in
$P_k^{r,r}(n-1)$.  If $(a_0,a_1,\dots,a_{n-1})$ is an incoming edge to this vertex then
$(a_1,a_2,\dots,a_{n})$ is an outgoing edge if and only if:
\[ a_n=\begin{cases}
0 \text{~or~} k/2 \text{~~if $k$ is even and $a_0=0$ or $a_0=k/2$}\\
a_0               \text{~~otherwise}
\end{cases} \]
\end{lemma}

\begin{proof}
Since both $(a_0,a_1,\dots,a_{n-1})$ and $(a_1,a_2,\dots,a_{n})$ have the same pseudoweight, it
follows that $f(a_0)=f(a_n)$, where $f$ is as defined in Definition~\ref{definition:pseudoweight}.
Hence, unless $k$ is even and $f(a_0)=f(a_n)=k/2$, this automatically implies $a_0=a_n$.
Conversely, if $f(a_0)=f(a_n)$ then $(a_0,a_1,\dots,a_{n-1})$ and $(a_1,a_2,\dots,a_{n})$ have the
same pseudoweight, and the result follows.
\end{proof}

The analysis for this function is somewhat more complex.  However, we do have a result similar to
Theorems~\ref{theorem:Hamming_characterisation}(i) and \ref{theorem:S_characteristic}.

\begin{theorem}  \label{theorem:P_characteristic_a}
Suppose $n\geq2$, $k>2$, and $n\leq r\leq n(k-1)$ (where $r\in\mathbb{Z}$ if $k$ is even and
$r=t/2$ for some $t\in\mathbb{Z}$ if $k$ is odd). Then $P_k^{r,r}(n-1)$ is Eulerian if and only if
\begin{itemize}
\item[(i)] $r\in\{n,(k-1)n\}$,
\item[(ii)] $r\in\{n+1,(k-1)n-1\}$ and $k\neq3$, or
\item[(iii)]  $r\in\{n-1+k/2,(k-1)(n-1)+k/2\}$ and $k$ odd.
\end{itemize}
\end{theorem}

\begin{proof}
Because of Lemma~\ref{lemma:P_balanced}, to show the Eulerian property holds it suffices to show
that the relevant subgraph is strongly connected.
\begin{itemize}
\item[(i)] $P_k^{n,n}(n-1)$ consists of a single edge, namely $1^n$ --- a loop from the vertex
    $1^{n-1}$ to itself; this is clearly Eulerian.  Similarly, the subgraph
    $P_k^{(k-1)n,(k-1)n}(n-1)$ consists of the single edge $(k-1)^n$ and thus is also Eulerian.
\item[(ii)] $P_4^{n+1,n+1}(n-1)$ contains all the $n$-tuples in the two necklaces $[1^{n-1},2]$
    and $[1^{n-1},0]$.  These two necklaces share a vertex, namely $1^{n-1}$, and hence the
    subgraph is strongly connected. If $k>4$, then $P_k^{n+1,n+1}(n-1)$ contains only a single
    necklace, namely $[1^{n-1},2]$, and thus the subgraph is strongly connected.  Analogous
    arguments hold for $P_4^{3n-1,3n-1}(n-1)$, which contains the necklaces $[3^{n-1},2]$ and
    $[3^{n-1},0]$, and $P_k^{n(k-1)-1,n(k-1)-1}(n-1)$, which contains the single necklace
    $[(k-1)^{n-1},k-2]$.
\item[(iii)] If $k$ is odd then $n-1+k/2$ is not an integer, and hence the only necklace in
    $P_k^{n-1+k/2,n-1+k/2}(n-1)$ is $[1^{n-1},0]$, and so $P_k^{n-1+k/2,n-1+k/2}(n-1)$ is
    Eulerian.  Similarly, $(k-1)(n-1)+k/2$ is not an integer, and hence the only necklace in
    $P_k^{(k-1)(n-1)+k/2,(k-1)(n-1)+k/2}(n-1)$ is $[(k-1)^{n-1},0]$ and so
    $P_k^{(k-1)(n-1)+k/2,(k-1)(n-1)+k/2}(n-1)$ is Eulerian.
\end{itemize}
It remains to show that in all other cases $P_k^{r,r}(n-1)$ is not Eulerian. First consider the
special case $k=3$ and $r\in\{n+1,(k-1)n-1\}$. Observe that $P_3^{n+1,n+1}$ contains two necklaces:
$[1^{n-1},2]$ and $[1^{n-2},0,0]$. These clearly do not share a vertex, and hence $P_3^{n+1,n+1}$
is not Eulerian.  Similarly, $P_3^{2n-1,2n-1}(n-1)$ contains the two necklaces. $[2^{n-1},1]$ and
$[2^{n-2},0,0]$. These clearly do not share a vertex, and hence $P_3^{2n-1,2n-1}(n-1)$ is not
Eulerian.

Next, suppose $k$ is odd, $n+1<r<(k-1)n-1$ and $r\notin\{n-1+k/2,(k-1)(n-1)+k/2\}$.  In this case
$P_k^{r,r}(n-1)$ will always contain more than one necklace. However, no two necklaces can share a
vertex since $w_p$ is one-to-one when $k$ is odd, and hence a vertex can only have at most one
incoming or outgoing edge with any given pseudoweight.  Hence, $P_k^{r,r}(n-1)$ is not strongly
connected.

Finally, suppose $k$ is even and $n+1< r< (k-1)n-1$.  Suppose $r=n+s(k-2)+t-1$, where $1\leq t<k-1$
and $0\leq s<n$.  Then there exists an edge in $P_k^{r,r}(n-1)$ of the form
$((k-1)^s,t,1^{n-s-1})$. No vertex in this necklace will have more than one incoming edge or
outgoing edge in $P_k^{r,r}(n-1)$ except in the case $t=k/2$.  If $t=k/2$ then the necklace
$[(k-1)^s,0,1^{n-s-1}]$ will share a vertex with the necklace $[(k-1)^s,t,1^{n-s-1}]$ but neither
necklace will share a vertex with any other edges in $P_k^{r,r}(n-1)$.

We next show that $P_k^{r,r}(n-1)$ will always contain another necklace not equal to either
$[(k-1)^s,t,1^{n-s-1}]$ (or $[(k-1)^s,0,1^{n-s-1}]$ if $t=k/2$), establishing that $P_k^{r,r}(n-1)$
is not strongly connected. We need to consider three cases.
\begin{itemize}
\item First suppose $s=0$, i.e.\ there exists an edge in $P_k^{r,r}(n-1)$ of the form
    $(t,1^{n-1})$, where $t>2$ since $r>n+1$.  Then there will also exist an edge of the form
    $(t-1,2,1^{n-2})$ in $P_k^{r,r}(n-1)$, since $n\geq2$, i.e.\ another necklace.
\item Next suppose $s>0$ and $t<k-2$; then, since $k\geq4$, there exists a further edge of the
    form $((k-1)^{s-1},k-2,t+1,1^{n-s-1})$ in $P_k^{r,r}(n-1)$.
\item Finally suppose $s>0$ and $t=k-2$.  Now, since $r< (k-1)n-1$, $n-s-1>0$.  Then the edge
    mentioned above takes the form $((k-1)^s,k-2,1^{n-s-1})$.  An additional edge is then
    $((k-1)^s,k-3,2,1^{n-s-2})$.
\end{itemize}
The result follows.
\end{proof}


The $k$ even case is somewhat simpler than the $k$ odd case, since if $k$ is even then
$w_p(\mathbf{a})\in\mathbb{Z}$ for any $k$-ary $n$-tuple $\mathbf{a}$.  We first need the following
preliminary lemma.

\begin{lemma}  \label{lemma:change_zeros_to_half_k}
Suppose $n\geq2$, $k>2$ is even, and $n\leq r\leq n(k-1)$. Suppose $\mathbf{a}$ is an edge (i.e.\ a
$k$-ary $n$-tuple) in $P_k^{r,r}(n-1)$ and $\bar{\mathbf{a}}$ is equal to $\mathbf{a}$ with every
$0$ changed to $k/2$. Then there exist directed paths in $P_k^{r,r}(n-1)$ from $\mathbf{a}$ to
$\bar{\mathbf{a}}$ and from $\bar{\mathbf{a}}$ to $\mathbf{a}$.
\end{lemma}

\begin{proof}
The result follows using repeated applications of Lemma~\ref{lemma:one_change}.  That is, to obtain
a directed path from $\mathbf{a}$ to $\bar{\mathbf{a}}$, simply use Lemma~\ref{lemma:one_change} to
obtain a directed path from $\mathbf{a}$ to an edge $\mathbf{a}'$ equal to $\mathbf{a}$ with a
single $0$ changed to $k/2$, and then a directed path from $\mathbf{a}'$ to an edge $\mathbf{a}''$
equal to $\mathbf{a}'$ with another $0$ changed to $k/2$, and so on until $\bar{\mathbf{a}}$ is
reached.  For the path in the opposite direction simply change the relevant entries of $k/2$ to $0$
one by one until $\mathbf{a}$ is reached.
\end{proof}

We can now give the main result.

\begin{theorem} \label{theorem:P_characteristic_k_even}
Suppose $n\geq2$ and $k>2$ is even. Then $P_k^{r,r+1}(n-1)$ is Eulerian for every $r\in \mathbb{Z}$
satisfying $n\leq r<(k-1)n$.
\end{theorem}

\begin{proof}
By Lemma~\ref{lemma:P_balanced} we need only show that $P_k^{r,r+1}(n-1)$ is strongly connected.
Also, by Lemma~\ref{lemma:change_zeros_to_half_k}, we need only show that there is a directed path
between any two edges in $P_k^{r,r+1}(n-1)$ whose corresponding $n$-tuples do not contain any
zeros.  But there is a simple graph isomorphism between the subgraph of $P_k^{r,r+1}(n-1)$ whose
edges are the $n$-tuples containing no zeros and $S_{k-1}^{r-n,r-n+1}(n-1)$ (by subtracting one
from every entry in an $n$-tuple), and so the result follows from
Theorem~\ref{theorem:S_characteristic}(ii).
\end{proof}

To analyse the $k$ odd case we first need the following.

\begin{lemma}  \label{lemma:k_odd_zero-free_path}
Suppose $n\geq2$, $k>2$ is odd, and $n+(k-1)/2\leq r \leq n(k-1)-(k-2)/2$, where $r=t/2$ for some
$t\in\mathbb{Z}$.  Then:
\begin{itemize}
\item[(i)] if $\mathbf{a}$ has pseudoweight $r$, there exists a directed path in
    $P_k^{r,r+0.5}(n-1)$ from $\mathbf{a}$ to $\mathbf{a}'$, where $\mathbf{a}'$ contains no
    entries equal to 0;
\item[(ii)] if $\mathbf{b}$ has pseudoweight $r+0.5$, there exists a directed path in
    $P_k^{r,r+0.5}(n-1)$ from $\mathbf{b}'$ to $\mathbf{b}$, where $\mathbf{b}'$ contains no
    entries equal to 0.
\end{itemize}
\end{lemma}

\begin{proof}
\begin{itemize}
\item[(i)]  First observe that, for any $k$-ary $n$-tuple $\mathbf{a}$ of pseudoweight $s$,
    replacing a single zero in $\mathbf{a}$ by $(k-1)/2$ results in in a $k$-ary $n$-tuple
    of pseudoweight $s-0.5$, and replacing a single zero by $(k+1)/2$ results in in a
    $k$-ary $n$-tuple of pseudoweight $s+0.5$. So, starting from $\mathbf{a}$ (of
    pseudoweight $r$), replace every $0$ one by one with alternately $(k+1)/2$ and
    $(k-1)/2$ until an $n$-tuple is reached (which we label $\mathbf{a}'$) containing no
    $0$ entries. In every case the pseudoweight of the result will be either $r$ or
    $r+0.5$, and hence by Lemma~\ref{lemma:one_change} there is a directed path from the
    starting $n$-tuple to the modified $n$-tuple in $P_k^{r,r+0.5}(n-1)$.  Joining these
    paths together yields the desired directed path from $\mathbf{a}$ to $\mathbf{a}'$.
\item[(ii)] The argument for (i) works in reverse, replacing every $0$ entry in
    $\mathbf{b}$ with alternately $(k-1)/2$ and $(k+1)/2$, yielding a directed path in
    $P_k^{r,r+0.5}(n-1)$ starting at an $n$-tuple $\mathbf{b}'$ containing no entries equal
    to $0$ and ending in $\mathbf{b}$.
\end{itemize}
\end{proof}

We can now give a companion result to Theorem~\ref{theorem:P_characteristic_k_even}.

\begin{theorem} \label{theorem:P_characteristic_k_odd}
Suppose $n\geq2$ and $k>2$ is odd. Then $P_k^{r,r+1}(n-1)$ is Eulerian for every $r\in \mathbb{Z}$
satisfying $n\leq r<(k-1)n$.
\end{theorem}

\begin{proof}
As previously, by Lemma~\ref{lemma:P_balanced} we need only show that $P_k^{r,r+1}(n-1)$ is
strongly connected.  Suppose $\mathbf{a}$ and $\mathbf{b}$ are distinct edges in
$P_k^{r,r+1}(n-1)$.

If $\mathbf{a}$ has no zeros then set $\mathbf{a}^*=\mathbf{a}$. Otherwise, if $\mathbf{a}$ has
pseudoweight $r$ or $r+0.5$, then, by Lemma~\ref{lemma:k_odd_zero-free_path}, there is a directed
path in $P_k^{r,r+1}(n-1)$ from $\mathbf{a}$ to an edge $\mathbf{a}^*$ which is zero-free.
Alternatively, if $\mathbf{a}$ has pseudoweight $r+1$ then let $\mathbf{a}'$ equal $\mathbf{a}$
with one zero replaced by $(k-1)/2$.  By Lemma~\ref{lemma:one_change}, since $\mathbf{a}'$ has
pseudoweight $r+0.5$, there is a directed path in $P_k^{r,r+1}(n-1)$ from $\mathbf{a}$ to
$\mathbf{a}'$. Moreover,  by Lemma~\ref{lemma:k_odd_zero-free_path} there is a directed path in
$P_k^{r,r+1}(n-1)$ from $\mathbf{a}'$ to an edge $\mathbf{a}^*$ which is zero-free.  That is, in
every case there is a directed path in $P_k^{r,r+1}(n-1)$ from $\mathbf{a}$ to an edge
$\mathbf{a}^*$ which is zero-free.

By a precisely analogous argument there is a directed path in $P_k^{r,r+1}(n-1)$ from an edge
$\mathbf{b}^*$ containing no zeros to $\mathbf{b}$.

But there exists a directed path in $P_k^{r,r+1}(n-1)$ from $\mathbf{a}^*$ to $\mathbf{b}^*$ using
exactly the same argument as employed in Theorem~\ref{theorem:P_characteristic_k_even} (observing
that $r$ is an integer).  The result follows.
\end{proof}

The separate results for $k$ odd and $k$ even (Theorems~\ref{theorem:P_characteristic_k_even} and
\ref{theorem:P_characteristic_k_odd}) can be combined to give the following, whose proof is
immediate using the same arguments as for Corollary~\ref{corollary:H_characterisation}.

\begin{corollary} \label{corollary:P_characteristic}
Suppose $n\geq2$, $k>2$ and $n\leq r< s\leq (k-1)n$, where $r,s\in\mathbb{Z}$. Then
$P_k^{r,s}(n-1)$ is Eulerian.
\end{corollary}

\section{Concluding remark}

It would be of interest to consider whether similar results could be obtained for other possible
permutation-agnostic weight functions.

%\bibliographystyle{amsplain}
%\bibliography{Coding}
\providecommand{\bysame}{\leavevmode\hbox to3em{\hrulefill}\thinspace}
\providecommand{\MR}{\relax\ifhmode\unskip\space\fi MR }
% \MRhref is called by the amsart/book/proc definition of \MR.
\providecommand{\MRhref}[2]{%
  \href{http://www.ams.org/mathscinet-getitem?mr=#1}{#2}
} \providecommand{\href}[2]{#2}
\begin{thebibliography}{10}

\bibitem{Alhakim22} A.~Alhakim, \emph{Designing preference functions for de {Bruijn} sequences
    with
  forbidden words}, Des. Codes Cryptogr. \textbf{90} (2022), no.~10,
  2319--2335.

\bibitem{Alhakim24a} A.~Alhakim, C.~J. Mitchell, J.~Szmidt, and P.~R. Wild, \emph{Orientable
  sequences over non-binary alphabets}, Cryptogr. Commun. \textbf{16} (2024),
  1309--1326.

\bibitem{Burns92} J.~Burns and C.~J. Mitchell, \emph{Coding schemes for two-dimensional
    position
  sensing}, Tech. Report HPL--92--19, January 1992,
  \url{https://www.chrismitchell.net/HPL-92-19.pdf}.

\bibitem{Burns93} \bysame, \emph{Coding schemes for two-dimensional position sensing},
  Cryptography and Coding III (M.~J. Ganley, ed.), Oxford University Press,
  1993, pp.~31--66.

\bibitem{Cameron25} B.~Cameron, A.~G{\"u}ndo{\u g}an, and J.~Sawada, \emph{Cut-down de {Bruijn}
  sequences}, Discrete Math. \textbf{348} (2025), 114024.

\bibitem{Chroman21} Z.~Chroman, M.~Kwan, and M.~Singhal, \emph{Lower bounds for superpatterns
    and
  universal sequences}, J. Comb. Theory {A} \textbf{182} (2021), 105467.

\bibitem{Chung92} F.~Chung, P.~Diaconis, and R.~Graham, \emph{Universal cycles for
    combinatorial
  structures}, Discrete Math. \textbf{{\bf 110}} (1992), 43--59.

\bibitem{Compeau11} P.~E.~C. Compeau, P.~A. Pevzner, and G.~Tesler, \emph{How to apply de
    {Bruijn}
  graphs to genome assembly}, Nat. Biotechnol. \textbf{29} (2011), 987--991.

\bibitem{Gabric20} D.~Gabric, J.~Sawada, A.~Williams, and D.~Wong, \emph{A successor rule
  framework for constructing $k$-ary de {Bruijn} sequences and universal
  cycles}, {IEEE} Trans. Inf. Theory \textbf{66} (2020), no.~1, 679--687.

\bibitem{Gibbons85} A.~Gibbons, \emph{Algorithmic graph theory}, Cambridge University Press,
  Cambridge, 1985.

\bibitem{Hollmann97} H.~D.~L. Hollmann and J.~H. {van Lint}, \emph{Universal sequences}, Appl.
  Algebra Eng. Commun. Comput. \textbf{8} (1997), no.~5, 347--352.

\bibitem{Jackson09} B.~Jackson, B.~Stevens, and G.~Hurlbert, \emph{Research problems on {Gray}
  codes and universal cycles}, Discrete Math. \textbf{309} (2009), no.~17,
  5341--5348.

\bibitem{Knuth97} D.~E. Knuth, \emph{The art of computer programming, {Volume} 1: {Fundamental}
  algorithms}, 3rd ed., Addison-Wesley, Reading, Mass., 1997.

\bibitem{Lempel85} A.~Lempel and M.~Cohn, \emph{Design of universal test sequences for {VLSI}},
  {IEEE} Trans. Inf. Theory \textbf{31} (1985), no.~1, 10--17.

\bibitem{Li23} M.~Li, Y.~Jiang, and D.~Lin, \emph{Properties of the cycles that contain all
  vectors of weight {$\le k$}}, Des. Codes Cryptogr. \textbf{91} (2023), no.~1,
  221--239.

\bibitem{Mitchell25a} C.~J. Mitchell and P.~R. Wild, \emph{Orientable and negative orientable
  sequences}, Discrete Appl. Math. \textbf{377} (2025), 242--259.

\bibitem{Mitchell26a} \bysame, \emph{Constructing orientable and negative orientable sequences
    with
  asymptotically optimal period}, 2026, Available at
  \url{https://arxiv.org/abs/2603.18646}.

\bibitem{Peng23} K.-Y. Peng, H.-S. Hsiao, and J.-Y. Chang, \emph{Length-adaptive linear
    position
  sensing system based on de-{Bruijn} sequence}, 2023 {IEEE} SENSORS, Vienna,
  Austria, October 29 -- Nov. 1, 2023, {IEEE}, 2023, pp.~1--4.

\bibitem{Penne10} R.~Penne, \emph{A note on certain de {Bruijn} sequences with forbidden
  subsequences}, Discrete Math. \textbf{310} (2010), no.~4, 966--969.

\bibitem{Petriu98} E.~M. Petriu, \emph{Absolute position measurement using pseudo-random binary
  encoding}, IEEE Instrum. Meas. Mag. \textbf{1} (1998), no.~3, 19--23.

\bibitem{Ruskey92} F.~Ruskey, C.~D. Savage, and T.~M.~Y. Wang, \emph{Generating necklaces}, J.
  Algorithms \textbf{13} (1992), no.~3, 414--430.

\bibitem{Ruskey12} F.~Ruskey, J.~Sawada, and A.~Williams, \emph{De {Bruijn} sequences for
  fixed-weight binary strings}, {SIAM} J. Discrete Math. \textbf{26} (2012),
  no.~2, 605--617.

\bibitem{Sawada11} J.~Sawada, B.~Stevens, and A.~Williams, \emph{De {Bruijn} sequences for the
  binary strings with maximum density}, {WALCOM:} Algorithms and Computation
  --- 5th International Workshop, {WALCOM} 2011, New Delhi, India, February
  18--20, 2011. Proceedings (N.~Katoh and A.~Kumar, eds.), Lecture Notes in
  Computer Science, vol. 6552, Springer, 2011, pp.~182--190.

\bibitem{Tan13} S.~Tan and J.~O. Shallit, \emph{Sets represented as the length-n factors of a
  word}, Combinatorics on Words --- 9th International Conference, {WORDS} 2013,
  Turku, Finland, September 16--20. Proceedings (J.~Karhum{\"{a}}ki,
  A.~Lepist{\"{o}}, and L.~Q. Zamboni, eds.), Lecture Notes in Computer
  Science, vol. 8079, Springer, 2013, pp.~250--261.

\bibitem{Yamaguchi05} K.~Yamaguchi and H.~Imai, \emph{Periodic sequences for absolute type
    shaft
  encoders}, Applied Algebra, Algebraic Algorithms and Error-Correcting Codes,
  8th International Conference, AAECC-8, Tokyo, Japan, August 20--24, 1990.
  Proceedings (S.~Sakata, ed.), Lecture Notes in Computer Science, vol. 508,
  Springer, Berlin, Heidelberg, 1991.

\end{thebibliography}

\end{document}
